	\begin{description}
	\item[Grundsätzlich:]Nachdem der Karate-D\={o} Ka die Grundschule und die vielfältige Formschule ausgiebig trainiert hat, beginnt mit der Wegschule eine neue Herausforderung für den Dan-Träger. Die Vorbereitungszeit ab der erfolgreich bestandenen Prüfung zum 1.\,Kyu beträgt mindestens 12 Monate. Mit Graduierung zum 1. Dan erfolgt auch die Einordnung in die Yudansha, die technischen Experten. Dan-Träger zu sein bedeutet stets Vorbild zu sein. 
	\item[Voraussetzung]zur richtigen Ausführung und Darstellung 
		\begin{enumerate}[nosep]
			\item Geistige Einstellung zur Kampfkunst
			\item Konditionelle körperliche Reife
			\item Schnelligkeit und Reaktionsfähigkeit
			\item Distanzgefühl
			\item Die richtige Wahl der Waffe
			\item Die Kenntnis der Atemi-Punkte
			\item Die richtige Entscheidung G\={o}jin Jitsu oder Jissen Jitsu
			\item Selbstvertrauen zur eigenen Technik
			\item Atmung und Kiai
			\item Dojo Kun
		\end{enumerate}
	\end{description}
	\renewcommand{\labelenumi}{\alph{enumi}.)}
	\renewcommand{\labelenumii}{\arabic*.}
	{\scriptsize 
	\begin{tabularx}{\textwidth}{lX}
		\textit{Vorausgesetzte Techniken}:& Zuki-Waza, Uke-Waza, Geri-Waza und Dachi-Waza der bis hierhin erlernten Kihon-Ido, Kata, Kata-Bunkai und Yakusoku-Kumite
	\end{tabularx}\\}
	\begin{enumerate}[nosep]
		\item \textit{\textbf{Kihon-Ido}}:
		\begin{enumerate}[nosep]
			\item Harai-Otoshi-Uke/Gyaku-Zuki J\={o}dan/Mae-Geri Ch\={u}dan\tabto{335pt}(Zenkutsu-Dachi)
			\item Age-Uke/Ren-Zuki J\={o}dan-Ch\={u}dan\tabto{335pt}(Sanchin-Dachi)
			\item Mawashi-Empi-Uke/Uraken-Uchi/Harai-Otoshi-Uke/Gyaku-Zuki Ch\={u}dan\tabto{335pt}(Shiko-Dachi 45°)
			\item Shuto-Uke/Mae-Geri Ch\={u}dan\tabto{335pt}(Neko-Ashi-Dachi)
		\end{enumerate}
		\vspace{12pt}
		\item \textit{\textbf{Kata}}:\tabto{200pt}\textit{\textbf{Kata\,Bunkai}}:
		\begin{enumerate}[nosep]
			\item Sanseru\tabto{200pt}Sanseru
		\end{enumerate}
		\vspace{12pt}
		\item \textit{\textbf{Partnerformen (2 aus 3)} - gilt für alle Dan-Prüfungen}:
		\begin{enumerate}[nosep]
			\item Yakusoku-Kumite:\tabto{200pt}3 Kumite-Ura\,\&\,2 Nage-Waza (frei auszuwählen)
			\item Shiai-Kumite:\tabto{200pt}2 Kämpfe nach DKV Reglement\\\tabto{200pt}\textit{entfällt für Sportler über 35 Jahre}
			\item Goshin-Jitsu-Kumite:\tabto{200pt}Abwehr \& Konter gegen Packen, Würgen, Halten\\\textit{(Jissen-Jitsu, Selbstverteidigung)}\tabto{200pt} von vorn, von der Seite und von hinten
		\end{enumerate}
		\vspace{12pt}
		\item \textit{\textbf{Weitere Verbandsvorgaben}}:
			\begin{enumerate}[nosep]
			\item Vorlage eines Kampfrichterlehrgangsnachweis
			\item Besuch eines Dansha-Lehrgangs der Stilrichtung G\={o}j\={u}-Ry\={u} Karate-D\={o}
			\item Unsere Abteilung empfiehlt die fortlaufende Teilnahme an Lehrgängen auf Landes-, oder Bundesebene, um den Horizont des eigenen Karate beständig zu erweitern
			\end{enumerate}
	\end{enumerate}